\documentclass[12pt]{article}
\usepackage{amsmath, amssymb}
\usepackage[a4paper, margin=1in]{geometry}

\begin{document}

\section*{Slide 1: Title Slide}
\textbf{CSE400 - Fundamentals of Probability in Computing} \\
\textbf{Lecture 21 and L22: Inequalities and Tail Bound} \\
Instructor: \textbf{Dhaval Patel, PhD} \\
Affiliation: SEAS, Ahmedabad University \\
Date: March 24 and 26, 2025  

\section*{Slide 2: Acknowledgement}
Acknowledgement to:  
\begin{itemize}
\item Mor Harchol-Balter, PhD  
\item J. Nelson Professor of Computer Science  
\item Carnegie Mellon University  
\end{itemize}

Reference Book:  
\begin{itemize}
\item \textit{Introduction to Probability for Computing}  
\item Publisher: Cambridge University Press (2024)  
\item ISBN: 978-1-009-30907-3  
\end{itemize}

\section*{Slide 3: Outline}
\begin{enumerate}
\item Tail and Motivation of Tail Bound  
\item Markov’s Inequality  
\begin{itemize}
\item Definition and Proof  
\item Example  
\end{itemize}
\item Chebyshev’s Inequality  
\begin{itemize}
\item Definition and Proof  
\item Example  
\end{itemize}
\item Chernoff Bound and Poisson Tail  
\begin{itemize}
\item Definition and Proof  
\item Example  
\end{itemize}
\item Jensen’s Inequality  
\begin{itemize}
\item Convex Functions  
\item Proof and Example  
\end{itemize}
\item Case Study and System-Level Understanding  
\end{enumerate}

\section*{Slide 4: Tails – Definition}
\textbf{Definition:}  
The tail of a random variable \( X \) is:
\[
P(X > x)
\]

\section*{Slide 5: Tails – Motivation}
\textbf{Why do we care about tails?}
\begin{itemize}
\item Jobs delayed more than 24 hours  
\item Packets dropped due to full buffer  
\end{itemize}

\section*{Slide 6: Tails – Key Idea}
\textbf{Key idea:}
\begin{itemize}
\item Mean $\rightarrow$ average behavior (easy)  
\item Tail $\rightarrow$ rare extreme events (hard)  
\end{itemize}

\section*{Slide 7: Tails – Difficulty}
\textbf{Why are tails hard to compute?}
\begin{itemize}
\item Requires summing many small probabilities  
\end{itemize}

\section*{Slide 8–10: Tails Example 1 (Binomial Case)}
\textbf{Question:}  
Suppose you are distributing \( n \) jobs among \( n \) servers at random.  
What is the probability that a particular server gets \( \geq k \) jobs?

\textbf{Intuition:}
\begin{itemize}
\item Jobs are assigned randomly  
\item Most servers get few jobs  
\item Occasionally, one server gets many jobs  
\end{itemize}

This is a \textbf{tail event (rare overload)}  

\textbf{Model:}
\[
X \sim \text{Binomial}(n, p)
\]

\textbf{Exact Tail Probability:}
\[
P(X \ge k) = \sum_{i=k}^{n} \binom{n}{i} p^i (1-p)^{n-i}
\]

\section*{Slide 11–13: Tails Example 2 (Poisson Case)}
\textbf{Question:}  
Jobs arrive according to a Poisson process with rate \( \lambda \).  
What is the probability that \( \geq k \) jobs arrive in one hour?

\textbf{Intuition:}
\begin{itemize}
\item Jobs arrive randomly over time  
\item Most hours have moderate arrivals  
\item Sometimes bursts occur  
\end{itemize}

This is a \textbf{tail event (sudden spike in load)}  

\textbf{Model:}
\[
X \sim \text{Poisson}(\lambda)
\]

\textbf{Exact Tail Probability:}
\[
P(X \ge k) = \sum_{i=k}^{\infty} e^{-\lambda} \frac{\lambda^i}{i!}
\]

\section*{Slide 14: Why Tail Bounds?}
Exact tail probabilities:
\[
P(X \ge k) = \sum_{i=k}^{\infty} e^{-\lambda} \frac{\lambda^i}{i!}
\]
\[
P(X \ge k) = \sum_{i=k}^{n} \binom{n}{i} p^i (1-p)^{n-i}
\]

\textbf{Observation:}
\begin{itemize}
\item These are complicated, long sums  
\end{itemize}

\section*{Slide 15: Idea of Tail Bounds}
Instead of computing exact value:
\begin{itemize}
\item Find an \textbf{upper bound}  
\item Easier to compute  
\item Still gives a guarantee  
\end{itemize}

\textbf{Goal:}
\[
\text{Exact: } P(X \ge k) = ?
\]
\[
\text{Bound: } P(X \ge k) \le \text{simple expression}
\]

\textbf{Key Idea:}  
Approximate safely instead of computing exactly  

\section*{Slide 16–17: Running Example Setup}
We will evaluate different bounds on the same example.

\textbf{Experiment:}
\begin{itemize}
\item Flip a fair coin \( n \) times  
\end{itemize}

\textbf{Distribution:}
\[
X = \text{number of heads}
\]
\[
X \sim \text{Binomial}(n, \tfrac{1}{2})
\]

\textbf{Mean:}
\[
E[X] = \frac{n}{2}
\]

\textbf{Question:}
\[
P\left(X \ge \frac{3n}{4}\right)
\]

\section*{Slide 18: Markov’s Inequality – Definition}
\textbf{Key Idea:}
\begin{itemize}
\item Uses only the mean  
\item Large values are rare  
\end{itemize}

\textbf{Markov Inequality:}
\[
P(X \ge a) \le \frac{E[X]}{a}
\]

\section*{Slide 19: Markov’s Inequality – Proof}
Consider:
\[
E[X] = \sum_{x=0}^{\infty} x \cdot p_X(x)
\]

Restrict to \( x \ge a \):
\[
E[X] \ge \sum_{x=a}^{\infty} x \cdot p_X(x)
\]

Replace \( x \ge a \) with \( a \):
\[
\ge \sum_{x=a}^{\infty} a \cdot p_X(x)
\]
\[
= a \sum_{x=a}^{\infty} p_X(x)
\]
\[
= a \cdot P(X \ge a)
\]

Thus:
\[
E[X] \ge a \cdot P(X \ge a)
\]
\[
P(X \ge a) \le \frac{E[X]}{a}
\]

\section*{Slide 20–21: Markov on Running Example}
\[
X \sim \text{Binomial}(n, \tfrac{1}{2}), \quad E[X] = \frac{n}{2}
\]

We want:
\[
P\left(X \ge \frac{3n}{4}\right)
\]

Apply Markov:
\[
P\left(X \ge \frac{3n}{4}\right) \le \frac{E[X]}{3n/4}
\]
\[
= \frac{n/2}{3n/4}
\]
\[
= \frac{2}{3}
\]

\textbf{Observation:}  
This is a poor bound (does not depend on \( n \))  

\section*{Slide 22–24: Markov Real-Life Example}
\textbf{Scenario:} Income distribution  
\[
E[X] = 10,000
\]

\textbf{Question:}  
Fraction earning \( \ge 50,000 \)

Apply Markov:
\[
P(X \ge 50,000) \le \frac{10,000}{50,000}
\]
\[
= 0.2
\]

\textbf{Result:}  
At most 20\% of people earn $\ge 50,000$  

\section*{Slide 25–29: Markov Solved Example 1}
\textbf{Problem:}  
Average grade = 70, A-grade cutoff = 90  

\subsection*{Step 1}
Let \( X \ge 0 \), \( E[X] = 70 \)

\subsection*{Step 2}
Markov inequality:
\[
P(X \ge a) \le \frac{E[X]}{a}
\]

\subsection*{Step 3}
Substitute:
\[
P(X \ge 90) \le \frac{70}{90}
\]

\subsection*{Step 4}
Simplify:
\[
= \frac{7}{9}
\]

\textbf{Result}
\[
P(X \ge 90) \le 0.778
\]

\section*{Slide 30–34: Markov Solved Example 2}
\textbf{Problem:}  
Average requests = 5000/sec  
Find \( P(X \ge 20,000) \)

\subsection*{Step 1}
\( E[X] = 5000 \)

\subsection*{Step 2}
\( a = 20,000 \)

\subsection*{Step 3}
\[
P(X \ge a) \le \frac{E[X]}{a}
\]

\subsection*{Step 4}
\[
P(X \ge 20,000) \le \frac{5000}{20,000}
\]
\[
= \frac{1}{4}
\]

\textbf{Result}
\[
P(X \ge 20,000) \le 0.25
\]

\section*{Slide 35: Markov Homework Problem}
Given:  
\( E[X] = 40,000 \)

(a) Bound \( P(X \ge 120,000) \)  

(b) Can Markov bound \( P(X < 10,000) \)?  

Hint:  
Markov applies to right tail \( P(X \ge a) \)  

\section*{Slide 36–38: Chebyshev’s Inequality}
\textbf{Key Idea:}
\begin{itemize}
\item Uses mean + variance  
\item Controls deviation from mean  
\end{itemize}

\textbf{Chebyshev Inequality:}
\[
P(|X - \mu| \ge a) \le \frac{\text{Var}(X)}{a^2}
\]

\textbf{Proof Idea:}
\[
P(|X - \mu| \ge a) = P((X - \mu)^2 \ge a^2)
\]
\[
\le \frac{E[(X - \mu)^2]}{a^2}
\]
\[
= \frac{\text{Var}(X)}{a^2}
\]

\section*{Slide 39: Chebyshev on Running Example}
\[
X \sim \text{Binomial}(n, \tfrac{1}{2})
\]
\[
\mu = \frac{n}{2}, \quad \text{Var}(X) = \frac{n}{4}
\]
\[
P\left(X \ge \frac{3n}{4}\right)
= P\left(X - \frac{n}{2} \ge \frac{n}{4}\right)
\]
\[
\le P\left(|X - \mu| \ge \frac{n}{4}\right)
\]
\[
\le \frac{n/4}{(n/4)^2}
= \frac{1}{n}
\]

\section*{Slide 40–42: Chebyshev Real-Life Example}
Given:
\[
\mu = 70, \quad \sigma = 10
\]

\[
P(X \le 40 \text{ or } X \ge 100)
= P(|X - 70| \ge 30)
\]

\[
\le \frac{10^2}{30^2}
= \frac{100}{900}
= \frac{1}{9}
\]

\textbf{Result:}  
At most \( \frac{1}{9} \) students are far from mean  

\section*{Slide 43–47: Chebyshev Solved Example 1}
Given:
\[
\mu = 70, \quad \sigma = 5 \Rightarrow \sigma^2 = 25
\]

\[
P(X \ge 90) \le P(|X - 70| \ge 20)
\]

\[
\le \frac{25}{400}
= \frac{1}{16}
\]

\textbf{Result}
\[
P(X \ge 90) \le 0.0625
\]

\section*{Slide 48–52: Chebyshev Solved Example 2}
Given:
\[
\mu = 50, \quad \sigma^2 = 16
\]

\[
X \notin [30, 70] \Rightarrow |X - 50| \ge 20
\]

\[
\le \frac{16}{400}
= \frac{1}{25}
\]

\textbf{Result}
\[
P(|X - 50| \ge 20) \le 0.04
\]

\section*{Slide 53: Chebyshev Homework}
Given:
\[
\mu = 100, \quad \sigma = 2
\]

Defective if:
\[
|X - 100| > 6
\]

\[
P(|X - \mu| \ge a) \le \frac{\sigma^2}{a^2}
\]

\section*{Slide 54–58: Chernoff Bound – Key Idea}
\begin{itemize}
\item Markov $\rightarrow$ uses \( X \)  
\item Chebyshev $\rightarrow$ uses \( (X - \mu)^2 \)  
\item Chernoff $\rightarrow$ uses \( e^{tX} \)  
\end{itemize}

\textbf{Insight:}
\begin{itemize}
\item Exponential grows very fast  
\item Large values amplified  
\item Tail events amplified  
\end{itemize}

\textbf{Better transformation ⇒ tighter bound}

\section*{Slide 59–61: Chernoff Core Trick}
\[
P(X \ge a) = P(e^{tX} \ge e^{ta})
\]

\[
P(e^{tX} \ge e^{ta}) \le \frac{E[e^{tX}]}{e^{ta}}
\]

\section*{Slide 62–64: Chernoff Bound (Final Form)}
\[
P(X \ge a) \le \min_{t > 0} \left\{ \frac{E[e^{tX}]}{e^{ta}} \right\}
\]

\textbf{Key Insight:}
\begin{itemize}
\item Holds for all \( t > 0 \)  
\item Choose best \( t \) → tightest bound  
\end{itemize}

\textbf{Much stronger than Markov \& Chebyshev}

\section*{Slide 65–68: Chernoff Lower Tail}
\[
P(X \le a) = P(e^{tX} \ge e^{ta})
\]

\[
P(X \le a) \le \min_{t<0} \left\{ \frac{E[e^{tX}]}{e^{ta}} \right\}
\]

\section*{Slide 69–70: Chernoff for Poisson Tail}
\[
X \sim \text{Poisson}(\lambda)
\]

\[
E[e^{tX}] = \sum_{i=0}^{\infty} e^{ti} \cdot \frac{e^{-\lambda}\lambda^i}{i!}
\]

\[
= e^{-\lambda} \sum_{i=0}^{\infty} \frac{(\lambda e^t)^i}{i!}
\]

\[
= e^{-\lambda} e^{\lambda e^t}
\]

\[
= e^{\lambda(e^t - 1)}
\]

\end{document}